
\documentclass[12pt,a4paper]{article}

\usepackage[utf8]{inputenc}
\usepackage[T1]{fontenc}
\usepackage[spanish]{babel}
\usepackage[margin=2.5cm]{geometry}
\usepackage{lmodern}
\usepackage{xcolor}
\usepackage{enumitem}
\usepackage{csquotes}
\usepackage[hidelinks]{hyperref}

\usepackage[
    backend=biber,
    style=ieee,
    sorting=none
]{biblatex}

\addbibresource{referencias.bib}

\definecolor{azul}{RGB}{35,55,80}

\setlength{\parindent}{0pt}
\setlength{\parskip}{7pt}

\newcommand{\ficha}[7]{

\subsection*{#1}

\noindent
{\color{azul}\rule{\linewidth}{0.8pt}}

\textbf{Referencia IEEE:}

\fullcite{#2}

\textbf{DOI verificado:}

\href{https://doi.org/#3}{#3}

\textbf{Resumen del artículo:}

#4

\textbf{Dato concreto:}

#5

\textbf{Nota personal:}

#6

\textbf{Uso en TA-IND-01:}

#7

\vspace{6mm}
}

\begin{document}

% ==============================
% PORTADA
% ==============================

\begin{titlepage}
\begin{center}

{\Large\bfseries
UNIVERSIDAD TÉCNICA ESTATAL
DE QUEVEDO}

\vspace{1.5cm}

{\large
FACULTAD DE CIENCIAS
DE LA INGENIERÍA}

\vspace{1.5cm}

{\Large\bfseries
FICHAS BIBLIOGRÁFICAS}

\vspace{0.5cm}

{\large
Fundamentos de sistemas distribuidos
y arquitecturas distribuidas}

\vspace{1.5cm}

\textbf{Actividad:}

FORM-FB-01

\vspace{0.5cm}

\textbf{Asignatura:}

Aplicaciones Distribuidas

\vspace{0.5cm}

\textbf{Estudiante:}

Emanuel Pino Juliana Romina

\vspace{0.5cm}

\textbf{Docente:}

Ing. Guerrero Ulloa Gleiston Ciceron

\vspace{0.5cm}

\textbf{Curso:}

Séptimo Software A

\vfill

Quevedo - Ecuador

Noviembre de 2026

\end{center}
\end{titlepage}

% ==============================
% INTRODUCCIÓN
% ==============================

\section{Introducción}

En esta actividad se revisaron
cinco artículos académicos
relacionados con los sistemas
distribuidos y los microservicios.

Los artículos tratan sobre
el mantenimiento de aplicaciones.
También explican cómo evolucionan
los microservicios y qué dificultades
pueden aparecer al utilizarlos.

El objetivo es conocer mejor
estos temas y reunir información
que pueda servir para el trabajo
individual TA-IND-01.

% ==============================
% CRITERIOS
% ==============================

\section{Criterios de selección}

Para realizar esta actividad
se seleccionaron cinco artículos
publicados entre 2023 y 2026.

Las fuentes fueron escogidas
porque están relacionadas con
los sistemas distribuidos
y las arquitecturas
de microservicios.

También se tomó en cuenta
que los artículos pertenecieran
a revistas académicas
y contaran con un DOI.

Para revisar cada fuente
se consideraron los siguientes
aspectos:

\begin{enumerate}[nosep]
\item El título del artículo.
\item Los autores y el año.
\item La revista donde fue publicado.
\item El problema que investiga.
\item Los resultados principales.
\item Un dato concreto del estudio.
\end{enumerate}

% ==============================
% FICHAS
% ==============================

\section{Fichas bibliográficas}

% ==============================
% FICHA 1
% ==============================

\ficha
{Ficha 1. Mantenimiento de microservicios}
{tighilt2023}
{10.1016/j.jss.2023.111755}
{
El artículo trata sobre los problemas
que pueden aparecer cuando se
desarrollan aplicaciones utilizando
microservicios. Aunque esta
arquitectura permite dividir un
sistema en partes independientes
también pueden presentarse errores
de diseño que dificultan su
mantenimiento. Los investigadores
proponen una herramienta llamada
MARS que permite identificar
diferentes antipatrones en los
microservicios. Para comprobar su
funcionamiento realizaron pruebas
con varios sistemas y analizaron
los resultados obtenidos.
El estudio muestra que detectar
estos problemas puede ayudar
a los desarrolladores a conocer
qué partes de una aplicación
necesitan mejoras. Esto resulta
útil para mantener los sistemas
organizados y facilitar
su mantenimiento.
}
{
El estudio presenta 16 antipatrones
y evalúa la herramienta MARS
en 24 sistemas.

Ubicación: resultados destacados
del artículo.
}
{
Este artículo me ayudó a entender
que trabajar con microservicios
no significa que una aplicación
esté bien diseñada automáticamente.
Pueden aparecer problemas que
después sean difíciles de corregir.
Me pareció interesante que existan
herramientas para identificar
esos errores.
}
{
Utilizaría esta información
para analizar si el sistema
elegido tiene una arquitectura
organizada y fácil de mantener.
}

% ==============================
% FICHA 2
% ==============================

\ficha
{Ficha 2. Evolución de microservicios}
{assuncao2023}
{10.1016/j.jss.2023.111788}
{
Esta investigación estudia cómo
cambian los microservicios durante
el desarrollo y mantenimiento
de una aplicación. Los autores
analizaron repositorios de
proyectos de código abierto
para conocer qué modificaciones
se realizan con el tiempo.
Para ello revisaron los cambios
registrados por los desarrolladores
e identificaron diferentes
características de su evolución.
El estudio permite comprender
que los microservicios necesitan
actualizaciones constantes
y que algunos cambios pueden
afectar a otros componentes
del sistema. También muestra
la importancia de mantener
una buena organización durante
el desarrollo. Esta información
ayuda a entender mejor cómo
evolucionan las aplicaciones
distribuidas y qué dificultades
pueden presentarse.
}
{
El estudio analizó 11 repositorios
de microservicios de código abierto.

Se revisaron 7.319 commits.

Ubicación: resultados destacados
del artículo.
}
{
Lo que comprendí es que los
microservicios no permanecen
iguales después de crearlos
porque necesitan cambios
y mejoras.

También entendí que modificar
un servicio puede afectar
al funcionamiento de otros.
}
{
Este artículo me serviría
para explicar cómo se mantiene
y evoluciona un sistema distribuido.
}

% ==============================
% FICHA 3
% ==============================

\ficha
{Ficha 3. Migración hacia microservicios}
{mohottige2025}
{10.1016/j.infsof.2025.107732}
{
El artículo analiza la
transformación de sistemas
tradicionales hacia arquitecturas
basadas en microservicios.
Este proceso puede resultar
complicado porque una aplicación
monolítica tiene sus componentes
relacionados y no siempre
es fácil separarlos.
Los investigadores realizaron
una revisión de estudios
anteriores para identificar
métodos y herramientas que
se utilizan durante la migración.
También analizaron los aspectos
que deben considerarse para
evaluar los resultados.
La investigación permite conocer
diferentes alternativas para
realizar este proceso y los
problemas que pueden surgir.
El trabajo sirve como referencia
para comprender qué se necesita
antes de cambiar la arquitectura
de una aplicación existente.
}
{
El artículo identifica métodos
y herramientas que se utilizan
para transformar sistemas
tradicionales en microservicios.

Ubicación: resumen del artículo.
Apartados Context y Objectives.
}
{
Este artículo me pareció útil
porque estoy aprendiendo sobre
la transformación de aplicaciones
tradicionales en sistemas web.

Entendí que cambiar una arquitectura
no consiste solamente en separar
el código. También es necesario
revisar cómo funcionan
sus componentes.
}
{
Utilizaría esta información
para comparar una arquitectura
monolítica con una arquitectura
basada en microservicios.
}

% ==============================
% FICHA 4
% ==============================

\ficha
{Ficha 4. Comunicación mediante Service Mesh}
{nicolas2024}
{10.1007/s10586-024-04342-5}
{
El artículo presenta una propuesta
para mejorar la integración de
patrones de procesamiento en
aplicaciones que utilizan
microservicios. Los autores
estudian las limitaciones
que pueden existir al administrar
la comunicación y el procesamiento
entre servicios. Para abordar
este problema proponen una solución
basada en service mesh que permite
organizar diferentes patrones
de procesamiento. La investigación
describe las etapas necesarias
para construir y desplegar estos
componentes y presenta una
evaluación de su funcionamiento.
Este trabajo ayuda a comprender
que además de desarrollar
microservicios independientes
también es necesario contar
con mecanismos que permitan
administrar sus conexiones
y facilitar su funcionamiento
dentro del sistema.
}
{
La propuesta considera tres fases:

\begin{itemize}[nosep]
\item Diseño y construcción.
\item Generación.
\item Despliegue.
\end{itemize}

Ubicación: sección 2 del artículo.
}
{
Antes de revisar este tema
no tenía muy claro para qué
servía un service mesh.

Ahora comprendo mejor que puede
ayudar a organizar la comunicación
entre microservicios.

Me parece importante porque
un sistema distribuido necesita
que sus componentes trabajen juntos.
}
{
Utilizaría esta información
para explicar cómo se comunican
los servicios y cómo se puede
administrar esa comunicación.
}

% ==============================
% FICHA 5
% ==============================

\ficha
{Ficha 5. Consistencia en sistemas distribuidos}
{bacao2026}
{10.1007/s00607-026-01701-5}
{
Este artículo estudia los
problemas de consistencia
que pueden aparecer en
aplicaciones desarrolladas
con microservicios orientados
a eventos. Cuando diferentes
servicios manejan información
relacionada pueden existir
dificultades para mantener
los datos y procesos actualizados.
Los autores realizaron una revisión
de investigaciones para identificar
los principales problemas
y las soluciones que se han
propuesto. El estudio organiza
la consistencia en tres aspectos:
datos y procesos y evolución
del sistema. También explica
que estos problemas deben
considerarse durante el diseño
de una arquitectura distribuida.
La investigación ayuda
a comprender la importancia
de coordinar los servicios
para mantener un funcionamiento
correcto.
}
{
El artículo identifica tres tipos
de consistencia:

\begin{itemize}[nosep]
\item Consistencia de datos.
\item Consistencia de procesos.
\item Consistencia de evolución.
\end{itemize}

Ubicación: secciones de resultados
y conclusiones.
}
{
Este artículo me ayudó a comprender
que uno de los problemas de los
sistemas distribuidos es mantener
la información correcta
entre varios servicios.

Me pareció interesante porque
no basta con que cada microservicio
funcione por separado.

También deben coordinarse
para evitar inconsistencias.
}
{
Utilizaría esta fuente
para analizar cómo un sistema
distribuido mantiene sus datos
actualizados y evita problemas
de consistencia.
}

% ==============================
% REFERENCIAS
% ==============================

\newpage

\nocite{
tighilt2023,
assuncao2023,
mohottige2025,
nicolas2024,
bacao2026
}

\printbibliography[
title={Referencias bibliográficas}
]

% ==============================
% DECLARACIÓN DE IA
% ==============================

\section*{Declaración de uso
de inteligencia artificial}

Para realizar esta actividad
utilicé ChatGPT como herramienta
de apoyo para organizar
las fuentes bibliográficas
y preparar el documento en LaTeX.

Los datos y resultados citados
deben comprobarse con
los artículos originales
antes de entregar el trabajo.

\end{document}

